\documentclass{article}

\usepackage{fancyhdr}
\usepackage{extramarks}
\usepackage{amsmath}
\usepackage{amsthm}
\usepackage{amsfonts}
\usepackage{tikz}
\usepackage[plain]{algorithm}
\usepackage{algpseudocode}
\usepackage{listings}
\usepackage{booktabs}
\usepackage{xcolor}
\usepackage[english]{babel}
\usepackage[T1]{fontenc}
\usepackage{lmodern,mathrsfs}
\usepackage{xparse}
\usepackage[inline,shortlabels]{enumitem}
\setlist{topsep=2pt,itemsep=2pt,parsep=0pt,partopsep=0pt}
\usepackage[dvipsnames]{xcolor}
\usepackage[utf8]{inputenc}
\usepackage[a4paper,top=0.5in,bottom=0.2in,left=0.5in,right=0.5in,footskip=0.3in,includefoot]{geometry}
\usepackage[most]{tcolorbox}
\tcbuselibrary{minted} % tcolorbox minted library, required to use the "minted" tcb listing engine (this library is not loaded by the option [most])
\usepackage{minted} % Allows input of raw code, such as Python code
% \usepackage[colorlinks]{hyperref}


\usetikzlibrary{automata,positioning}

\tcbset{
    pythoncodebox/.style={
        enhanced jigsaw,breakable,
        colback=gray!10,colframe=gray!20!black,
        boxrule=1pt,top=2pt,bottom=2pt,left=2pt,right=2pt,
        sharp corners,before skip=10pt,after skip=10pt,
        attach boxed title to top left,
        boxed title style={empty,
            top=0pt,bottom=0pt,left=2pt,right=2pt,
            interior code={\fill[fill=tcbcolframe] (frame.south west)
                --([yshift=-4pt]frame.north west)
                to[out=90,in=180] ([xshift=4pt]frame.north west)
                --([xshift=-8pt]frame.north east)
                to[out=0,in=180] ([xshift=16pt]frame.south east)
                --cycle;
            }
        },
        title={#1}, % Argument of pythoncodebox specifies the title
        fonttitle=\sffamily\bfseries
    },
    pythoncodebox/.default={}, % Default is No title
    %%% Starred version has no frame %%%
    pythoncodebox*/.style={
        enhanced jigsaw,breakable,
        colback=gray!10,coltitle=gray!20!black,colbacktitle=tcbcolback,
        frame hidden,
        top=2pt,bottom=2pt,left=2pt,right=2pt,
        sharp corners,before skip=10pt,after skip=10pt,
        attach boxed title to top text left={yshift=-1mm},
        boxed title style={empty,
            top=0pt,bottom=0pt,left=2pt,right=2pt,
            interior code={\fill[fill=tcbcolback] (interior.south west)
                --([yshift=-4pt]interior.north west)
                to[out=90,in=180] ([xshift=4pt]interior.north west)
                --([xshift=-8pt]interior.north east)
                to[out=0,in=180] ([xshift=16pt]interior.south east)
                --cycle;
            }
        },
        title={#1}, % Argument of pythoncodebox specifies the title
        fonttitle=\sffamily\bfseries
    },
    pythoncodebox*/.default={}, % Default is No title
}

% Custom tcolorbox for Python code (not the code itself, just the box it appears in)
\newtcolorbox{pythonbox}[1][]{pythoncodebox=#1}
\newtcolorbox{pythonbox*}[1][]{pythoncodebox*=#1} % Starred version has no frame

% Basic Document Settings
\topmargin=-0.45in
\evensidemargin=0in
\oddsidemargin=0in
\textwidth=6.5in
\textheight=9.0in
\headsep=0.25in
\linespread{1.1}

\pagestyle{fancy}
\lhead{\hmwkAuthorName}
\chead{\hmwkClass\ (\hmwkClassInstructor): \hmwkTitle}
\rhead{\firstxmark}
\lfoot{\lastxmark}
\cfoot{\thepage}
\renewcommand\headrulewidth{0.4pt}
\renewcommand\footrulewidth{0.4pt}
\setlength\parindent{0pt}

% Homework Details
\newcommand{\hmwkTitle}{Asignación 2}
\newcommand{\hmwkDueDate}{December 14, 2025}
\newcommand{\hmwkClass}{ESMA 6661}
\newcommand{\hmwkClassInstructor}{Israel Almodovar}
\newcommand{\hmwkAuthorName}{\textbf{Alejandro Ouslan}}

% Title Page
\title{
	\vspace{2in}
	\textmd{\textbf{\hmwkClass:\ \hmwkTitle}}\\
	\normalsize\vspace{0.1in}\small{Due\ on\ \hmwkDueDate}\\
	\vspace{0.1in}\large{\textit{\hmwkClassInstructor}}
	\vspace{3in}
}

\author{\hmwkAuthorName}
\date{}


% Begin document
\begin{document}
\maketitle
\pagebreak
\pagebreak

\section*{Problem 1: Definitions}

\begin{enumerate}[(a)]
	\item \textbf{Sufficiency principle:}
	\item \textbf{Factorization theorem:}
	\item \textbf{Likelihood Function:}
	\item \textbf{Minimal sufficient:}
	\item \textbf{Complete statistics:}
\end{enumerate}

\section*{Problem 2: TBD}

In each of these exercises, assume that the random variables $X_1,\ldots,X_n$ form a random
sample of size $n$ from the distribution specified in that exercise, and show that the statistic
$T$ specified in the exercise is a sufficient statistic for the parameter.

\begin{enumerate}[(a)]
	\item The Bernoulli distribution with parameter $p$, which is unknown $(0<p<1); T= \sum_{n}^{i=1}X_i$.
	\item The geometric distribution with parameter $p$, which is unknown $(0<p<1); T= \sum_{n}^{i=1}X_i$.
	\item The negative binomial distribution with parameters $r$ and $p$, where $r$ is known and $p$
	      is unknown $(0<p<1); T = \sum_{n}^{i = 1}X_i$.
	\item The normal distribution for which the mean $\mu$ is known and the variance $\sigma^2>0$ is
	      unknown; $T = \sum_{n}^{i=1}(X_i - \mu)^2$.
	\item The gamma distribution with parameters $\alpha$ and $\beta$, where the value of $\alpha$ is known and
	      the value of $\beta$ is unknown $(\beta>0);T= \bar{X}_n$.
	\item The gamma distribution with parameters $\alpha$ and $\beta$, where the value of $\beta$ is known and
	      the value of $\beta$ is unknown $(\beta>0);T= \prod^{n}_{i=1}X_n$.
	\item The beta distribution with parameters $\alpha$ and $\beta$, where the value of $\beta$ is known and
	      the value of $\beta$ is unknown $(\beta>0);T= \prod^{n}_{i=1}X_n$.


\end{enumerate}

\end{document}
